\chapter{Backend service — nghiệp vụ và API}

\section{Tổ chức mã nguồn}

\begin{longtable}{L{4.0cm} L{9.6cm}}
\toprule
\textbf{Thư mục} & \textbf{Nội dung}\\
\midrule
\endhead
\code{app/main.py} & Điểm vào ứng dụng: đăng ký \en{middleware}, bộ định tuyến, bộ xử lý ngoại lệ, vòng đời khởi động/tắt.\\
\code{app/core/} & Cấu hình (\code{config.py}), động cơ cơ sở dữ liệu bất đồng bộ, bảo mật (băm mật khẩu, JWT), xác thực Firebase, phụ thuộc FastAPI, \en{middleware}, Redis, bộ nhớ đệm, Celery.\\
\code{app/models/} & 25 tệp mô hình ORM SQLAlchemy, khoảng 70 bảng.\\
\code{app/schemas/} & Lược đồ Pydantic v2 — hợp đồng đầu vào/ra của API.\\
\code{app/routes/} & 41 mô-đun định tuyến.\\
\code{app/services/} & 45 dịch vụ nghiệp vụ (không chứa mã HTTP).\\
\code{app/crud/} & Lớp truy cập dữ liệu tái sử dụng cho 8 miền lớn.\\
\code{app/clients/} & \code{ai\_service\_client.py} — cửa duy nhất gọi sang AI service.\\
\code{app/tasks/} & 9 mô-đun tác vụ nền Celery.\\
\code{alembic/} & Di trú lược đồ cơ sở dữ liệu.\\
\code{scripts/} & Gieo dữ liệu, kiểm toán nội dung, sinh bài tập bằng AI, sao lưu.\\
\bottomrule
\caption{Cấu trúc thư mục backend-service}
\end{longtable}

\section{Chuỗi middleware}

\en{Middleware} là lớp bọc quanh ứng dụng: yêu cầu đi vào phải xuyên qua từ ngoài vào
trong, phản hồi đi ra xuyên qua theo chiều ngược lại. Thứ tự đăng ký quyết định thứ tự
chạy — lớp đăng ký sau cùng nằm ngoài cùng.

\begin{center}
\code{PrivateNetworkAccess $\rightarrow$ CORS $\rightarrow$ RateLimit $\rightarrow$ RequestID $\rightarrow$ RequestLogging $\rightarrow$ App}
\end{center}

\begin{longtable}{L{4.6cm} L{9.0cm}}
\toprule
\textbf{Lớp} & \textbf{Chức năng}\\
\midrule
\endhead
\code{TrustedHostMiddleware} & Chỉ chấp nhận yêu cầu có tiêu đề \code{Host} nằm trong danh sách cho phép; chỉ bật ở môi trường sản xuất.\\
\code{PrivateNetworkAccess} & Gắn tiêu đề CORS-RFC1918 để trình duyệt Chrome cho phép truy cập mạng nội bộ trong lúc phát triển.\\
\code{CORSMiddleware} & Kiểm soát nguồn gốc chéo (\en{cross-origin}) — chỉ cho phép miền \code{lexilingo.me} và các nguồn cấu hình.\\
\code{RateLimitMiddleware} & Cửa sổ trượt trên Redis; mặc định 300 yêu cầu/phút khi phát triển và 120 khi sản xuất, trần 5.000/giờ.\\
\code{RequestIDMiddleware} & Sinh \code{X-Request-ID} để lần vết một yêu cầu qua toàn bộ nhật ký. Tác vụ Celery cũng nhận mã tương tự dạng \code{task:\{id\}}.\\
\code{RequestLoggingMiddleware} & Ghi phương thức, đường dẫn, mã trạng thái và độ trễ.\\
\code{limit\_request\_body} & Từ chối thân yêu cầu vượt \code{MAX\_REQUEST\_BODY\_BYTES}.\\
\bottomrule
\caption{Chuỗi middleware của backend}
\end{longtable}

\section{Nhóm định tuyến}

Toàn bộ API dùng tiền tố \code{/api/v1}. Bốn mươi mốt mô-đun được nhóm theo miền:

\begin{longtable}{L{3.9cm} L{9.7cm}}
\toprule
\textbf{Nhóm} & \textbf{Mô-đun và chức năng}\\
\midrule
\endhead
Xác thực \& người dùng & \code{auth} (đăng ký, đăng nhập, làm mới thẻ, Google/Facebook qua Firebase, quên/đổi mật khẩu), \code{users} (hồ sơ, tuỳ chọn, ảnh đại diện), \code{devices} (đăng ký thiết bị nhận thông báo đẩy), \code{rbac} (vai trò, quyền), \code{entitlements} (quyền lợi gói trả phí), \code{referral} (giới thiệu bạn bè).\\
Nội dung học & \code{courses}, \code{course\_categories}, \code{learning} (phiên học, bắt đầu/hoàn thành bài), \code{vocabulary} (ngân hàng từ, ôn tập), \code{concepts} (khái niệm tri thức), \code{mistakes} (sổ tay lỗi sai).\\
Tiến độ \& năng lực & \code{progress} (tiến độ, chuỗi ngày, hoạt động hằng ngày), \code{proficiency} (đánh giá CEFR sáu kỹ năng, cổng thi thăng bậc), \code{learner\_state} (điểm vào nội bộ cho quan sát từ AI service).\\
Động lực học & \code{gamification} (thành tích, ví, cửa hàng, kho đồ, bảng xếp hạng), \code{xp} (giao dịch điểm kinh nghiệm), \code{challenges} (thử thách), \code{games} (trò chơi mini — mô-đun lớn nhất).\\
Nội dung thực tế & \code{news}, \code{podcasts}, \code{books}, \code{youtube} — nội dung ngoài đời thực kèm câu hỏi/quiz.\\
Thông báo & \code{notifications}, \code{reminders} (tuỳ chọn nhắc học), \code{notification\_campaign} (chiến dịch gửi hàng loạt).\\
Tác tử AI & \code{content\_agent} (sinh khoá học), \code{ranking\_agent} (xếp hạng, giải đấu), \code{admin\_ai\_proxy} (uỷ nhiệm gọi AI service từ trang quản trị), \code{ai\_audit} (nhật ký quyết định AI).\\
Quản trị \& vận hành & \code{admin\_courses}, \code{admin\_gamification}, \code{admin\_system}, \code{user\_management}, \code{analytics}, \code{monitoring}, \code{product\_events}, \code{integrations}, \code{health}, \code{well\_known}.\\
\bottomrule
\caption{Bốn mươi mốt mô-đun định tuyến theo miền}
\end{longtable}

\section{Xác thực và phân quyền}

\subsection{Luồng xác thực}

Hệ thống hỗ trợ ba đường đăng nhập, tất cả cùng quy về một loại thẻ nội bộ:

\begin{enumerate}[leftmargin=1.4em]
  \item \textbf{Email/mật khẩu} — mật khẩu băm bằng \code{bcrypt} gọi trực tiếp
  (\code{bcrypt.hashpw}), không qua \code{passlib} vì thư viện này có lỗi tương thích với
  bcrypt 4.x.
  \item \textbf{Google / Facebook} — client lấy thẻ Firebase, backend xác minh bằng
  Firebase Admin SDK rồi phát hành JWT của chính hệ thống.
  \item \textbf{Đối tác} — khoá API lưu ở bảng \code{partner\_api\_keys}.
\end{enumerate}

Thẻ truy cập (\en{access token}) có tuổi thọ ngắn, thẻ làm mới (\en{refresh token}) lưu ở
bảng \code{refresh\_tokens} và có thể thu hồi. Ứng dụng Flutter chủ động làm mới
\textbf{trước khi} thẻ hết hạn thay vì đợi lỗi 401 — thay đổi này (commit
\code{557d29c4}) vừa giảm độ trễ vừa cắt chi phí gọi Firebase.

\subsection{Phân quyền theo vai trò}

Bốn bảng \code{roles}, \code{permissions}, \code{role\_permissions} và \code{audit\_logs}
tạo nên hệ RBAC đầy đủ. Quyền được kiểm tra bằng \en{dependency} của FastAPI, tức là khai
báo ngay trên chữ ký hàm định tuyến chứ không nằm rải rác trong thân hàm. Mọi thao tác
quản trị nhạy cảm ghi lại vào \code{audit\_logs}.

\section{Các dịch vụ nghiệp vụ chính}

\begin{longtable}{L{5.0cm} L{8.6cm}}
\toprule
\textbf{Dịch vụ} & \textbf{Trách nhiệm}\\
\midrule
\endhead
\code{proficiency\_service} & Chấm điểm sáu kỹ năng CEFR, tính điểm tổng hợp có trọng số, quyết định điều kiện thăng bậc. Xem Chương~9.\\
\code{learner\_state} & Tiến hoá trạng thái khái niệm theo FSRS + luật delta; điểm vào duy nhất \code{record\_concept\_observation}.\\
\code{learner\_state\_outbox} & Mẫu \en{outbox}: ghi sự kiện quan sát vào cùng giao dịch rồi mới xác nhận, tránh mất dữ liệu khi mạng đứt.\\
\code{fsrs\_scheduler\_service}, \code{game\_fsrs\_service} & Lập lịch ôn tập cho từ vựng và cho trò chơi.\\
\code{xp\_service} & Cấp điểm kinh nghiệm, áp trần theo ngày và theo nguồn, hệ số chuỗi ngày, chống cấp trùng.\\
\code{achievement\_checker\_service} & Kiểm tra và mở khoá thành tích theo \en{trigger} (sự kiện kích hoạt) thay vì quét toàn bộ.\\
\code{streak\_service} & Chuỗi ngày học liên tiếp, khiên bảo vệ chuỗi.\\
\code{rank\_service}, \code{ranking\_agent} & Xếp hạng, giải đấu theo tuần, đặt lại bảng xếp hạng.\\
\code{game\_scoring\_service} & Chấm điểm trò chơi, quy đổi sang kỹ năng CEFR tương ứng.\\
\code{content\_agent\_apply} & Đường duy nhất dựng khoá học $\rightarrow$ chương $\rightarrow$ bài học kèm bài tập từ kết quả của tác tử nội dung.\\
\code{content\_agent\_validation} & Làm sạch và kiểm tra hình dạng bài tập trước khi ghi.\\
\code{vocabulary\_catalog} & Danh mục từ vựng và chính sách chọn từ.\\
\code{skill\_history\_service} & Tổng hợp thống kê kỹ năng theo ngày; cắt tỉa nhật ký chi tiết.\\
\code{entitlement\_service}, \code{revenuecat\_client} & Quyền lợi gói trả phí, đồng bộ với RevenueCat.\\
\code{push\_notification\_service}, \code{reminder\_service} & Thông báo đẩy và nhắc học theo lịch FSRS.\\
\code{user\_deletion\_service} & Xoá tài khoản và dữ liệu liên quan theo yêu cầu quyền riêng tư.\\
\code{quota\_manager}, \code{api\_cache\_service} & Hạn ngạch gọi API theo người dùng; bộ nhớ đệm phản hồi trên Redis.\\
\bottomrule
\caption{Dịch vụ nghiệp vụ backend}
\end{longtable}

\section{Bất biến nội dung khoá học}

Một bài học chỉ \textbf{chơi được} khi \code{content.exercises} không rỗng. Bất biến này
được thực thi ở nhiều tầng, vì mỗi tầng từng bị vi phạm:

\begin{itemize}[leftmargin=1.4em]
  \item Cột \code{Lesson.exercise\_count} được đồng bộ tự động bằng \en{hook}
  \code{@validates("content")}; không được gán tay. Mọi câu hỏi "bài này học được chưa?"
  phải dùng cột này, không dùng \code{total\_exercises}.
  \item Các điểm cuối hướng người học (lộ trình, chi tiết khoá học) ẩn bài học rỗng;
  \code{/learning/lessons/\{id\}/start} và \code{/content} trả lỗi 409 ở môi trường sản xuất.
  \item Xuất bản khoá học bị chặn nếu còn bài học rỗng
  (\code{CourseCRUD.publish\_blockers}), chỉ kiểm tra tại thời điểm chuyển
  bản nháp $\rightarrow$ đã xuất bản.
  \item Mọi chỉnh sửa bài học từ trang quản trị đều đi qua \code{\_apply\_lesson\_update}
  — dùng chung bởi \code{PUT /lessons/\{id\}} và \code{PUT /lessons/\{id\}/content}.
\end{itemize}

\begin{luuy}[Hình dạng bài tập không được kiểm tra khi ghi]
Hệ thống kiểm tra "có bài tập hay không" nhưng không kiểm tra "bài tập có đúng hình dạng
không". Ba dạng đã từng hỏng: dạng ghép cặp phải nối bằng dấu \code{", "} vì bộ chấm tách
theo dấu phẩy; dạng \code{arrange\_the\_sentence} có đáp án \textbf{không} nằm trong danh
sách lựa chọn nên không dùng được phép kiểm tra chung; dạng \code{image\_based\_choice}
không có đường ống ảnh nên tuyệt đối không được sinh ra.
\end{luuy}

\section{Tác vụ nền Celery}

Celery chạy hai tiến trình riêng: \code{worker} (thực thi) và \code{beat} (hẹn giờ), dùng
Redis làm môi giới. Cấu hình đặt \code{task\_acks\_late=True} và
\code{worker\_prefetch\_multiplier=1} để tác vụ chỉ được xác nhận sau khi hoàn thành.

\begin{longtable}{L{5.6cm} L{3.2cm} L{4.6cm}}
\toprule
\textbf{Tác vụ} & \textbf{Lịch} & \textbf{Mục đích}\\
\midrule
\endhead
\code{scan\_fsrs\_reminders} & theo chu kỳ cấu hình & Tìm khái niệm đến hạn ôn, đẩy nhắc nhở.\\
\code{send\_streak\_alerts} & 20:00 hằng ngày & Cảnh báo sắp mất chuỗi ngày học.\\
\code{send\_word\_of\_day} & 08:00 hằng ngày & Từ vựng trong ngày.\\
\code{cleanup\_expired\_content\_agent\_uploads} & 03:15 hằng ngày & Dọn tệp tải lên hết hạn.\\
\code{cleanup\_learner\_observations} & 02:30 hằng ngày & Dọn quan sát đã áp dụng.\\
\code{prune\_exercise\_attempts} & 02:45 hằng ngày & Cắt nhật ký trả lời cũ hơn 90 ngày.\\
\code{snapshot\_skill\_scores} & 03:00 thứ Hai & Chụp ảnh điểm kỹ năng để tính xu hướng 7/30 ngày.\\
\code{auto\_league\_reset} & 00:05 thứ Hai & Đặt lại giải đấu tuần.\\
\code{prefetch\_news} & mỗi 6 giờ & Nạp trước tin tức.\\
\code{prefetch\_youtube} & mỗi 12 giờ & Nạp trước video.\\
\code{prefetch\_podcasts} & 00:00 hằng ngày & Nạp trước podcast.\\
\bottomrule
\caption{Lịch tác vụ nền}
\end{longtable}

\begin{luuy}[Vì sao có \code{snapshot\_skill\_scores}]
Cột \code{UserSkillScore.trend} đọc \code{score\_7d\_ago} và \code{score\_30d\_ago} —
hai cột mà trước đó \textbf{không có tiến trình nào ghi vào}. Xu hướng vì thế luôn bằng
không. Tác vụ tuần này là cách rẻ nhất để có lịch sử mà không cần thêm bảng lịch sử riêng.
\end{luuy}

\section{Xử lý lỗi}

Ba bộ xử lý ngoại lệ được đăng ký ở \code{main.py}, đảm bảo mọi lỗi trả về cùng một hình
dạng JSON:

\begin{lstlisting}[language=Python]
{
  "error": {
    "code": "VALIDATION_ERROR",
    "message": "...",
    "details": [ ... ]
  }
}
\end{lstlisting}

Nhờ vậy client chỉ cần một bộ phân tích lỗi duy nhất, và mã lỗi (\code{code}) là hằng số
ổn định để so khớp thay vì phải đọc thông điệp tiếng Anh.
